\section{Conclusion}
\label{sec:conclusion}

We test whether instruct LLMs favor early answer positions and whether this inflates multiple-choice scores, using a within-question rotation design on seven open-weight models and three benchmarks.
Answer order changes what these models answer, but not in the way the hypothesis states.
There is no shared preference for early positions: each model has its own stable pattern, and most do worse than their own average at A on \mmlu.
Because benchmark keys are close to uniform, the net effect on published scores is about one point or less.
The effect is concentrated on questions where the model is close to 50/50, is larger with five options than four because of the last slot, is largest for word-answer true/false in small models, and depends strongly on the label symbols.

The key takeaway is that, for these models, answer position makes multiple-choice scores unstable at the question level rather than inflated in aggregate.

Future work should, in order of value: rerun the same items on hosted frontier models with a generated-text read-out; separate label from position with a crossed design in which labels are permuted independently of slots; use full random shuffles on self-contained questions only, to remove the order-dependent-option confound; test whether the residual advantage of the published \mmlu order appears on a benchmark released after the models' training cut-offs; and repeat the true/false experiment with plain factual statements and several phrasings.
